\documentclass[11pt]{article}
\usepackage[T1]{fontenc}
\usepackage{lmodern}
\usepackage[margin=0.93in]{geometry}
\usepackage{amsmath,amssymb,amsthm,mathtools}
\usepackage{microtype,booktabs,array,longtable}
\usepackage{xcolor,hyperref,fancyhdr,enumitem}
\definecolor{ink}{HTML}{183149}
\definecolor{accent}{HTML}{176E86}
\hypersetup{colorlinks=true,linkcolor=accent,citecolor=accent,urlcolor=accent,
 pdftitle={RA-14: Bounds and a spectral-to-PCA reduction},
 pdfauthor={Sidney Holden},pdfsubject={Partial research result; not a complete solution to RA-14}}
\pagestyle{fancy}
\fancyhf{}
\fancyhead[L]{\small\textcolor{ink}{RA-14 \;|\; Spectral rank-$k$ approximation}}
\fancyhead[R]{\small Research note}
\fancyfoot[L]{\small\textcolor{ink}{Partial result --- September 12, 2026}}
\fancyfoot[R]{\thepage}
\setlength{\headheight}{14pt}
\setlength{\parskip}{4pt}
\setlength{\parindent}{0pt}
\setlength{\emergencystretch}{2em}
\numberwithin{equation}{section}
\newtheorem{theorem}{Theorem}[section]
\newtheorem{lemma}[theorem]{Lemma}
\newtheorem{proposition}[theorem]{Proposition}
\newtheorem{corollary}[theorem]{Corollary}
\theoremstyle{definition}
\newtheorem{definition}[theorem]{Definition}
\newtheorem{example}[theorem]{Example}
\newtheorem{remark}[theorem]{Remark}
\newcommand{\R}{\mathbb R}
\newcommand{\E}{\mathbb E}
\newcommand{\Prb}{\mathbb P}
\newcommand{\Id}{I}
\newcommand{\norm}[1]{\left\lVert#1\right\rVert}
\newcommand{\tr}{\operatorname{tr}}
\newcommand{\rank}{\operatorname{rank}}
\newcommand{\range}{\operatorname{range}}
\newcommand{\diag}{\operatorname{diag}}
\newcommand{\qsp}{q_{\mathrm{sp}}}
\newcommand{\eps}{\varepsilon}
\newcommand{\St}{\operatorname{Stief}}
\newcommand{\poly}{\operatorname{poly}}
\newcommand{\K}{\mathcal K}
\newcommand{\cgood}{\mathcal E}
\newcommand{\code}[1]{\texttt{#1}}

\begin{document}
\begin{titlepage}
\vspace*{0.45in}
{\large\sffamily\color{accent} OPEN PROBLEMS IN NUMERICAL LINEAR ALGEBRA}\par
\vspace{0.3in}
{\Huge\bfseries\color{ink} RA-14}\par
\vspace{0.12in}
{\LARGE\bfseries Bounds and a spectral-to-PCA reduction}\par
\vspace{0.2in}
{\large Exact two-sided matrix--vector queries}\par
\vspace{0.38in}
\begingroup
\setlength{\fboxsep}{12pt}
\noindent\colorbox{ink!6}{\parbox{\dimexpr\textwidth-24pt\relax}{
\textbf{Status: partial research result, not a complete solution.}\par
The proofs below give a universal rank lower bound, an exact padding reduction,
and matching bounds in a sufficiently-large-dimension regime that allows the
rank to grow. They do not characterize the query complexity for every
simultaneous choice of dimension, rank, and accuracy.
}}
\endgroup
\vspace{0.3in}
\textbf{Main reduction.} On a symmetric input with positive leading eigenvalues
and a suitable gap, an already valid spectral rank-$k$ approximation can be
converted into a strong PCA approximation with at most
\[
 k\left\lceil\frac{10}{\sqrt\eps}\right\rceil
\]
additional matrix--vector products. A weighted graph inequality and a
Fej\'er-type polynomial remove a logarithmic accuracy loss from a simpler
subspace-angle argument.
\vfill
\textbf{Sidney Holden}\par
{\small Biological Transport Networks, Center for Computational Biology\\
Flatiron Institute, Simons Foundation}\par
September 12, 2026\par
\smallskip
\small Generated mathematical research note. The reduction proofs are written
out; the two cited complexity theorems are imported. No claim of priority,
independent peer review, or machine-checked proof is made. The accompanying
code is a numerical illustration, not an oracle lower-bound certificate.
\end{titlepage}

\tableofcontents
\newpage

\section{Problem, scope, and results}
\label{sec:scope}

For $n\ge2$, $1\le k<n$, and $0<\eps<1/2$, an unknown $A\in\R^{n\times n}$
can be queried through either $Ax$ or $A^Tx$. Each vector product costs one
query. Queries may be fully adaptive; other computation is free. The required
output is $Z\in\R^{n\times k}$ with $Z^TZ=I_k$ and
\begin{equation}
 \Prb\!\left[\norm{A(I-ZZ^T)}_2\le(1+\eps)\sigma_{k+1}(A)\right]\ge0.99
 \quad\text{for every }A.
 \label{eq:target}
\end{equation}
The minimum worst-case query budget is $\qsp(n,k,\eps)$. This is the
RA-14 formulation in the supplied repository \cite{RA14}. Algorithms are
understood to be measurable, so the probabilities in the model are defined.

All logarithms below are natural. Universal constants are independent of
$n,k,\eps$. A block product with $b$ columns always costs $b$ vector queries.

\begin{theorem}[Universal bounds and a large-rank regime]
\label{thm:global}
For a universal constant $C_{\rm up}>0$,
\begin{equation}
 k\ \le\ \qsp(n,k,\eps)\ \le\
 \min\left\{n,\left\lceil C_{\rm up}\frac{k\log n}{\sqrt\eps}\right\rceil\right\}.
 \label{eq:global}
\end{equation}
Consequently $\qsp(n,k,\eps)=\Theta(n)$ whenever $k\ge n/2$, uniformly in
$\eps$. On the promised subclass $\rank(A)\le k$, exactly $k$ queries suffice
and are necessary.
\end{theorem}

The lower bound and promised-subclass statement are proved in
Section~\ref{sec:rank}. The $n$-query upper bound simply learns all columns.
The Krylov upper bound is the spectral guarantee of Musco--Musco
\cite[Algorithm 2 and Theorem 10]{MM15}, applied to $A^T$; the query accounting
is given in Section~\ref{sec:upper}.

\begin{theorem}[Matching bounds with polynomially large dimension]
\label{thm:large}
There exist universal constants $C_*,D,c_*>0$ such that
\begin{equation}
 n\ge C_*\left(\frac{k}{\eps}\right)^D
 \quad\Longrightarrow\quad
 c_*\frac{k\log n}{\sqrt\eps}
 \ \le\ \qsp(n,k,\eps)
 \ \le\ \left\lceil C_{\rm up}\frac{k\log n}{\sqrt\eps}\right\rceil.
 \label{eq:largedim}
\end{equation}
In particular, this statement is not restricted to constant $k$.
\end{theorem}

The lower bound in Theorem~\ref{thm:large} is a deduction from the published
PCA lower bound of Simchowitz--El Alaoui--Recht \cite{SAR18}, using the
postprocessing theorem proved in Section~\ref{sec:filter}. The exponent $D$
is not optimized or evaluated explicitly here. The dimension restriction is
essential to what has actually been proved in this note.

A further, rank-independent lower bound follows by exact padding:
\begin{equation}
 \qsp(n,k,\eps)\ge c_{\rm BN}\frac{\log(n-k+1)}{\sqrt\eps}
 \quad\text{if }n-k+1\ge C_{\rm BN}\eps^{-2.01},
 \label{eq:paddedBN}
\end{equation}
with universal constants and an enlarged universal threshold if necessary.
This transfers the bounded hard instances of Bakshi--Narayanan
\cite[Theorems 1.1, 5.2, and Section 6]{BN23}; the transfer itself is proved in
Section~\ref{sec:padding}.

\subsection*{What is not established}
Equations~\eqref{eq:global}, \eqref{eq:largedim}, and \eqref{eq:paddedBN} do
not match in all remaining parameter regimes. In particular, the proof does
not justify placing an $n$-cap around the high-dimensional lower bound and
then asserting it for every $n,k,\eps$. The upper bound permits that cap;
the lower-bound construction still has a dimension hypothesis. This is the
specific missing step, rather than a claim that the requested problem is
impossible to solve.

\section{A universal lower bound of \texorpdfstring{$k$}{k} queries}
\label{sec:rank}

The zero-error requirement on rank-$k$ inputs supplies a useful lower bound,
but a proof must handle both oracle directions and arbitrary adaptivity.
A random orthogonal projector alone is not an adequate construction when
$k$ is close to $n$: one can learn its small complementary subspace instead.
We use a nonsymmetric rank-$k$ family.

\begin{lemma}[Unrevealed Gaussian block]
\label{lem:gaussian}
Let $H\in\R^{n\times k}$ have independent standard Gaussian entries.
A deterministic adaptive procedure observes products $Hc$ and $H^Tx$.
Let $C\subset\R^k$ and $X\subset\R^n$ be the spans of the respective query
vectors, with dimensions $r$ and $s$. Conditional on the transcript,
\begin{equation}
 H=\bar H+P_{X^\perp}GP_{C^\perp},\qquad
 \bar H=P_XH+HP_C-P_XHP_C,
 \label{eq:gaussblock}
\end{equation}
where $\bar H$ is known and the remaining block is a fresh standard Gaussian
map from $C^\perp$ to $X^\perp$. Equivalently, $G$ can be taken to have fresh
independent standard Gaussian entries in fixed ambient coordinates.
\end{lemma}

\begin{proof}
Initially the statement is the definition of $H$. Conditional on a transcript,
the next query vector is fixed. A new column query depends, in the unrevealed
block, only on its component in $C^\perp$. If that component is nonzero,
choose it as the first vector of an orthonormal basis of $C^\perp$.
The query reveals one column of an independent Gaussian matrix, leaving
all other columns independent and unchanged. A row query has the same
argument with the first row of the remaining block. A redundant query
reveals nothing. These steps prove the assertion by induction, including
adaptive choices of the next vector. The two known restrictions are $HP_C$
and $P_XH$; inclusion--exclusion gives the displayed $\bar H$.
\end{proof}

\begin{lemma}[Independence of the row responses]
\label{lem:rowrank}
Before $k$ total nonredundant queries have been made,
$\rank(H^TX)=\dim X=s$ almost surely.
\end{lemma}

\begin{proof}
Consider a new row query $x\notin X$. Conditional on the preceding transcript,
$H^Tx$ has a nondegenerate Gaussian component with full support in
$C^\perp$, of dimension $k-r$. Since $r+s<k$, this dimension exceeds $s$.
The affine support of the new answer cannot be contained in the span of the
old $s$ row answers. Its probability of lying in that span is therefore zero.
The rank increases by one. Column queries do not alter the already obtained
row responses. Induction proves the claim.
\end{proof}

\begin{proposition}[Adaptive two-sided rank lower bound]
\label{prop:ranklower}
No algorithm making fewer than $k$ vector queries can satisfy
\eqref{eq:target} on every rank-$k$ input.
\end{proposition}

\begin{proof}
Set
\begin{equation}
 A=\begin{bmatrix}H^T\\0_{(n-k)\times n}\end{bmatrix}.
 \label{eq:rankfamily}
\end{equation}
Almost surely, $\rank A=k$. An $A$ query reveals $H^Tx$; an $A^T$ query
reveals $Hc$, where $c$ consists of the first $k$ coordinates of its input.
Thus these are exactly the queries in Lemma~\ref{lem:gaussian}.
Because $\sigma_{k+1}(A)=0$, success requires
\[
 \range Z=\range H.
\]
Fix the algorithm's random seed, and condition on its transcript after
$q<k$ queries. Write $S=\range Z$, $N=X^\perp$. We have $r+s\le q<k$.

If $N\not\subset S$, choose a nonzero $c\in C^\perp$. By
\eqref{eq:gaussblock}, $Hc$ has a Gaussian distribution with affine support
parallel to $N$. A proper intersection of that support with $S$ has
probability zero. Hence $\Prb[\range H\subset S\mid\text{transcript}]=0$.

If $N\subset S$, Lemma~\ref{lem:rowrank} gives
$P_X\range H=X$. Consequently $\range H+N=\R^n$.
The inclusions $\range H\subset S$ and $N\subset S$ would imply
$S=\R^n$, contradicting $\dim S=k<n$. Success is again impossible.

Thus success has probability zero for this input distribution after fixing
any seed. Integrating over the seed gives the same conclusion for randomized
algorithms. An algorithm succeeding with probability $0.99$ on every input
would also succeed with that probability under this distribution, a
contradiction.
\end{proof}

\begin{corollary}[Exact complexity under a rank promise]
\label{cor:promise}
On inputs promised to satisfy $\rank A\le k$, the query complexity is $k$.
\end{corollary}
\begin{proof}
Take $k$ independent Gaussian vectors $g_j$ and query $A^Tg_j$.
Their span equals the row space of $A$ almost surely. Orthonormalize and,
if the actual rank is less than $k$, extend to a $k$-dimensional subspace.
The residual is zero. Proposition~\ref{prop:ranklower} supplies necessity.
\end{proof}

\section{Upper bounds and vector-query accounting}
\label{sec:upper}

Querying $Ae_1,\ldots,Ae_n$ determines $A$. A singular-value decomposition
then gives a right singular basis satisfying the target with equality.
This proves the exact $n$-query cap without any spectral assumptions.

For the faster upper bound, the right-sided version of block Krylov uses
\[
 \range\!\left[A^TG,(A^TA)A^TG,\ldots,(A^TA)^tA^TG\right],
 \qquad G\in\R^{n\times k}\ \text{Gaussian}.
\]
The cited spectral theorem yields the $0.99$ guarantee for
$t=O(\log n/\sqrt\eps)$, with the exact-column branch used when it is
cheaper \cite{MM15}. A reorthogonalized implementation starts with $A^TG$
and alternates products by $A$ and $A^T$ on each newly produced block.
It also retains the images $AQ$ needed to form the small Ritz problem.
At depth $t$ this costs at most
\begin{equation}
 k(2t+2)
 \label{eq:uppercount}
\end{equation}
vector products, including compression. No block is charged as a single
vector product. Selecting the top right singular vectors of the already
known $AQ$ requires no further oracle access.

The underlying polynomial analysis can be written using
$\log(n/\eps)$. In the branch where $k\log n/\sqrt\eps<n$, one has
$\log(1/\eps)\le2\log n$ after a universal constant adjustment.
Outside that branch, learning all columns gives the desired cap.
Thus the displayed global upper bound is genuinely finite-dimensional.

Combining this cap with Proposition~\ref{prop:ranklower} gives
$n/2\le\qsp\le n$ for $k\ge n/2$, proving the large-rank assertion of
Theorem~\ref{thm:global}.

\section{What a valid spectral approximation says about angles}
\label{sec:angles}

\begin{lemma}[Uniform overlap]
\label{lem:angle}
Let $V\in\R^{n\times k}$ be a top right singular basis of $A$, and let
$a=\sigma_k(A)>0$. If $Z^TZ=I_k$ and
$\norm{A(I-ZZ^T)}_2\le\tau<a$, then
\begin{equation}
 \sigma_{\min}(V^TZ)^2\ge1-\frac{\tau^2}{a^2}.
 \label{eq:angle}
\end{equation}
\end{lemma}
\begin{proof}
For the leading left singular basis $U$ and diagonal singular block
$\Sigma$, one has $V^T=\Sigma^{-1}U^TA$. Therefore
$\norm{V^T(I-ZZ^T)}_2\le\tau/a$. Moreover,
\[
 V^TZZ^TV=I_k-V^T(I-ZZ^T)V
 \succeq\left(1-\frac{\tau^2}{a^2}\right)I_k.
\]
The eigenvalues on the left are the squared singular values of $V^TZ$.
\end{proof}

\begin{example}[A sharp obstruction to direct PCA inference]
\label{ex:sharp}
Take $n=2k+1$, $a=1+2\eps$, $b=1$, $\tau=1+\eps$, and
\[
 A=\diag(aI_k,0_k,b),\qquad
 Z=\begin{bmatrix}\alpha I_k\\\sqrt{1-\alpha^2}I_k\\0\end{bmatrix},
 \quad \alpha^2=1-\left(\frac\tau a\right)^2.
\]
Then $Z^TZ=I_k$, $\sigma_{k+1}(A)=b$, and
\begin{equation}
 \norm{A(I-ZZ^T)}_2=\max\{b,a\sqrt{1-\alpha^2}\}=\tau.
 \label{eq:sharpresidual}
\end{equation}
Thus the target is met exactly and the angle bound is sharp. However,
\[
 \frac{\tr(Z^TAZ)}{\sum_{i=1}^k\lambda_i(A)}=\alpha^2
 =\frac{2\eps+3\eps^2}{1+4\eps+4\eps^2}\longrightarrow0.
\]
A spectral approximation need not itself capture nearly all of the leading
PCA objective. A PCA lower bound therefore cannot simply be applied to the
original output $Z$ without a further argument.
\end{example}

\section{Spectral approximation to PCA in \texorpdfstring{$O(k/\sqrt\eps)$}{O(k/sqrt(epsilon))} extra queries}
\label{sec:filter}

This section proves the main reduction. The starting block is already a
valid spectral approximation; the postprocessor does not have to discover
such a block from an arbitrary random initialization.

\begin{theorem}[Deterministic warm-start postprocessing]
\label{thm:warm}
Let $M\in\R^{n\times n}$ be symmetric. Its top $k$ eigenvalues are positive,
$\lambda_i\ge a$, and all remaining eigenvalues $\mu_j$ satisfy
$|\mu_j|\le b$. Suppose $0<\eps<1/2$,
\begin{equation}
 a\ge(1+\tfrac32\eps)b,
 \qquad Z^TZ=I_k,
 \qquad \norm{M(I-ZZ^T)}_2\le(1+\eps)b.
 \label{eq:warmassumptions}
\end{equation}
There is a deterministic postprocessor, requiring neither $a$ nor $b$, which
uses at most $k\lceil10/\sqrt\eps\rceil$ further products with $M$ and returns
$\widehat V^T\widehat V=I_k$ with
\begin{equation}
 \tr(\widehat V^TM\widehat V)
 \ge\left(1-\frac\eps{200}\right)\sum_{i=1}^k\lambda_i.
 \label{eq:warmconclusion}
\end{equation}
The same lower bound holds with $M$ replaced by $|M|$ in the objective.
\end{theorem}

\subsection{A weighted graph inequality}
Assume $b>0$ and set $\tau=(1+\eps)b$. Work in an eigenbasis in which
$M=\diag(\Lambda_1,\Lambda_2)$, where $\Lambda_1=\diag(\lambda_i)$ and
$\Lambda_2=\diag(\mu_j)$. Since $a>\tau$, Lemma~\ref{lem:angle} says that
the leading block of $Z$ is invertible. Its range is therefore the graph
\[
 \range Z=\range\begin{bmatrix}I_k\\F_0\end{bmatrix}
\]
for a matrix $F_0\in\R^{(n-k)\times k}$. A vector perpendicular to this graph
has the form $(-F_0^Ty,y)$. Applying the residual bound to these vectors gives
\begin{equation}
 F_0(\Lambda_1^2-\tau^2I)F_0^T\preceq\tau^2I-\Lambda_2^2.
 \label{eq:weightedgraph}
\end{equation}
Define the positive definite diagonal matrices
$D_1=\Lambda_1^2-\tau^2I$ and $D_2=\tau^2I-\Lambda_2^2$.
Equation~\eqref{eq:weightedgraph} is equivalent to
\begin{equation}
 F_0=D_2^{1/2}KD_1^{-1/2},\qquad \norm K_2\le1,
 \qquad\norm K_F^2\le k.
 \label{eq:contract}
\end{equation}
Unlike an unweighted angle bound, this controls the amount of leakage into
each tail eigenvalue according to its distance from the residual threshold.

\subsection{A polynomial and its elementary bounds}
For an integer $m\ge1$, define the polynomial
\begin{equation}
 p_m(u)=\frac{T_m(u)-1}{u-1},\qquad p_m(1)=m^2,
 \label{eq:fejer}
\end{equation}
where $T_m(\cos\theta)=\cos(m\theta)$ is the Chebyshev polynomial.
The quotient has degree $m-1$. The trigonometric identity
\begin{equation}
 p_m(u)=m+2\sum_{j=1}^{m-1}(m-j)T_j(u)
 \label{eq:fejerexpansion}
\end{equation}
implies that $p_m$ is positive and nondecreasing for $u\ge1$.
For $-1\le u\le1$ it also gives
\begin{equation}
 0\le p_m(u)\le\min\left\{m^2,\frac2{1-u}\right\},
 \label{eq:fejertail}
\end{equation}
with the second bound omitted at $u=1$. Indeed,
$p_m(\cos\theta)=\bigl(\sin(m\theta/2)/\sin(\theta/2)\bigr)^2$.
The first bound follows by expressing the sine quotient as a sum of $m$
unit complex numbers, and the second from $|T_m(u)-1|\le2$.
For $u>1$,
\begin{equation}
 p_m(u)=\frac{\cosh(m\operatorname{arcosh}u)-1}{u-1}.
 \label{eq:fejerhead}
\end{equation}

\subsection{Trace loss after filtering}
Apply $p_m(M/b)$ to the starting block. Its range is the graph of
\begin{equation}
 F=p_m(\Lambda_2/b)F_0p_m(\Lambda_1/b)^{-1}.
 \label{eq:filteredgraph}
\end{equation}
An orthonormal basis of this graph is
$Y=[I;F](I+F^TF)^{-1/2}$. Put $C=(I+F^TF)^{-1}$. The trace loss satisfies
\begin{align}
 \sum_i\lambda_i-\tr(Y^TMY)
 &=\tr\bigl((\Lambda_1-aI)F^TFC\bigr)
   +\tr\bigl(F^T(aI-\Lambda_2)FC\bigr)\notag\\
 &\le\sum_{i,j}(\lambda_i-\mu_j)|F_{ji}|^2.
 \label{eq:traceloss}
\end{align}
For the first inequality use $0\preceq F^TFC\preceq F^TF$ and
$\Lambda_1-aI\succeq0$; for the second term use $0\prec C\preceq I$ and
$F^T(aI-\Lambda_2)F\succeq0$. No commutation of $\Lambda_1$ and $F^TF$ is
assumed.

Substituting \eqref{eq:contract} and \eqref{eq:filteredgraph} into
\eqref{eq:traceloss}, and using $\norm K_F^2\le k$, bounds the loss by $k$
times the largest scalar coefficient
\begin{equation}
 \mathcal B(\lambda,\mu)=
 \frac{(\lambda-\mu)(\tau^2-\mu^2)}{\lambda^2-\tau^2}
 \left(\frac{p_m(\mu/b)}{p_m(\lambda/b)}\right)^2,
 \quad\lambda\ge a,\quad |\mu|\le b.
 \label{eq:coefficient}
\end{equation}
For fixed $\mu$, $(\lambda-\mu)/(\lambda^2-\tau^2)$ is decreasing for
$\lambda>\tau$, because the numerator of its derivative is
\[
 \mu^2-\tau^2-(\lambda-\mu)^2<0.
\]
The denominator polynomial is nondecreasing and positive. It is therefore
safe to replace $\lambda$ by $a_0=(1+\frac32\eps)b$ when bounding
\eqref{eq:coefficient} from above.

\subsection{An explicit constant bound}
Choose
\begin{equation}
 m=\left\lceil\frac{10}{\sqrt\eps}\right\rceil.
 \label{eq:mchoice}
\end{equation}
Scale $b=1$ for the moment, write $e=\eps$, $s=1-\mu\in[0,2]$, and put
$a_0=1+\frac32e$, $t_0=1+e$. The numerator in \eqref{eq:coefficient} is
\[
 N=p_m(\mu)^2(\tfrac32e+s)(e+s)(2+e-s).
\]
If $s\ge e$, inequality~\eqref{eq:fejertail} gives
\[
 N\le\frac4{s^2}\cdot\frac52s\cdot2s\cdot\frac52=50.
\]
If $s\le e$, then $m\sqrt e\le10+\sqrt e<11$, so
\[
 N\le\frac{25}{2}m^4e^2
 <\frac{25}{2}\,11^4<200000.
\]
Thus $N<200000$ in both cases.

For $0\le x\le1$, the power series gives $\cosh x\le1+x^2$.
It follows that $\operatorname{arcosh}(1+\frac32e)\ge\sqrt e$.
Also $a_0^2-t_0^2=e+\frac54e^2\ge e$. Consequently,
\begin{equation}
 p_m(a_0)^2(a_0^2-t_0^2)
 \ge\frac{(\cosh10-1)^2}{(9/4)e}
 >\frac{4\cdot10^8}{9e}.
 \label{eq:denominator}
\end{equation}
The elementary numerical inequality used here has the exact certificate
\[
 \cosh10-1>
 \sum_{j=1}^{7}\frac{10^{2j}}{(2j)!}
 =\frac{5860808350}{567567}>10000.
\]
Combining numerator and denominator, then restoring $b$, proves
\begin{equation}
 \mathcal B(\lambda,\mu)<\frac{9}{2000}eb
 <\frac{eb}{200}.
 \label{eq:scalarfinal}
\end{equation}
Hence the filtered graph satisfies
\[
 \sum_i\lambda_i-\tr(Y^TMY)
 \le\frac{k\eps b}{200}
 \le\frac\eps{200}\sum_i\lambda_i.
\]

\subsection{Implementing the proof without spectral side information}
The vector block $p_m(M/b)Z$ lies in
\begin{equation}
 \K_{m-1}(M,Z)=\range[Z,MZ,\ldots,M^{m-1}Z].
 \label{eq:warmkrylov}
\end{equation}
The algorithm builds this space and returns its top $k$ algebraic Ritz
vectors. The variational principle ensures that their trace is at least
$\tr(Y^TMY)$. The algorithm does not evaluate $p_m$ or use the unknown $b$;
the polynomial only proves that a sufficiently good subspace is present.

To count products, orthogonalize successive frontier blocks. Each frontier
has at most $k$ columns. Compute its image under $M$, retain that image for
the compression, and use its component outside the current space as the
next frontier. Process frontiers numbered $0$ through $m-1$, including the
last image block needed for the Ritz matrix. At most $km$ products are used.
If the space becomes invariant or fills $\R^n$, stop earlier. All needed
entries of $Q^TMQ$ are already known.

Finally, $|M|-M\succeq0$, giving the claim for $|M|$. If $b=0$, the residual
in \eqref{eq:warmassumptions} is zero and $Z$ already spans the leading
space; the same algorithm can stop after detecting invariance. This proves
Theorem~\ref{thm:warm}.

\section{The growing-rank lower bound in large dimension}
\label{sec:SAR}

\subsection{The imported theorem, with its dimension condition retained}
We use the following consequence of \cite[Theorem 1, Proposition 2.1,
and Theorem 2.2]{SAR18}. For $0<g\le1/2$, set
\begin{equation}
 \lambda=\frac{1+\sqrt{g(2-g)}}{1-g},\qquad
 L=\lambda+\lambda^{-1},\qquad L-2=gL,
 \qquad M=W+\lambda UU^T,
 \label{eq:SARmodel}
\end{equation}
where $U$ is uniform in $\St(n,k)$ and $W$ is independent GOE, with
off-diagonal variance $1/n$ and diagonal variance $2/n$. Define
\begin{equation}
 \cgood_\gamma=
 \{\norm W_2+(1-\gamma)gL\le\lambda_k(M)\}
 \cap\{\lambda_1(M)\le(1+\gamma)L\}.
 \label{eq:goodevent}
\end{equation}
There are universal $c_0,c_2>0$ and polynomial dimension thresholds such that,
under the law conditioned on $\cgood_{1/2}$, every fully adaptive algorithm
using $T\le c_0k\log n/\sqrt g$ has probability at most $\exp(-n^{c_2})$ of
returning an orthonormal $\widehat V$ satisfying
\begin{equation}
 \tr(\widehat V^T|M|\widehat V)
 \ge(1-g/12)\sum_{i=1}^k\sigma_i(M).
 \label{eq:SARtarget}
\end{equation}
For fixed $\gamma=1/20$ and failure probability $0.01$, a further polynomial
threshold gives $\Prb(\cgood_{1/20})\ge0.99$. These results are imported,
not re-proved here.

\subsection{Checking every hypothesis of the postprocessor}
Set $g=\min\{4\eps,1/2\}$, so $\eps\le g\le4\eps$. On
$\cgood_{1/20}$, put $a=\lambda_k(M)$, $b=\sigma_{k+1}(M)$, and
$w=\norm W_2$. Interlacing for the positive rank-$k$ perturbation gives
$-w\le\lambda_j(M)\le w$ for all $j>k$. Also $a>w$ on this event, so the
leading $k$ eigenvalues are positive and are exactly the leading singular
values. Furthermore,
\[
 a-b\ge\frac{19}{20}gL,\qquad
 a\le\lambda_1(M)\le\frac{21}{20}L,
\]
and consequently
\begin{equation}
 \frac ba\le1-\frac{19}{21}g.
 \label{eq:SARgap}
\end{equation}
If $b=0$, the postprocessor's conclusion is immediate. Otherwise, when
$\eps\le1/8$,
\[
 \frac ab\ge\frac1{1-76\eps/21}
 \ge1+\frac{76}{21}\eps>1+\frac32\eps.
\]
When $\eps\ge1/8$,
\[
 \frac ab\ge\frac{42}{23}>\frac74\ge1+\frac32\eps.
\]
Thus any successful output of an RA-14 algorithm satisfies all hypotheses of
Theorem~\ref{thm:warm}. Its postprocessed output meets
\eqref{eq:SARtarget}, since $\eps/200<g/12$.

\subsection{Probability and query budget}
Suppose an RA-14 algorithm uses at most $q$ queries. On symmetric $M$, either
oracle direction is just a product with $M$, so the same algorithm is
admissible for the imported theorem. Append the deterministic postprocessor.
The total budget is at most
\begin{equation}
 T=q+k\left\lceil\frac{10}{\sqrt\eps}\right\rceil.
 \label{eq:composedbudget}
\end{equation}
Since $\cgood_{1/20}\subset\cgood_{1/2}$,
\[
 \Prb(\cgood_{1/20}\mid\cgood_{1/2})\ge0.99.
\]
The original approximation algorithm succeeds with probability at least
$0.99$ for every fixed $M$, even when we condition on a particular generating
pair $(W,U)$. Therefore the composed PCA success probability under the
conditioned hard law is at least $0.99^2=0.9801$.

For large enough dimension, this contradicts the imported theorem if the
right side of \eqref{eq:composedbudget} is at most
$c_0k\log n/\sqrt g$. Hence
\begin{equation}
 q>c_0\frac{k\log n}{\sqrt g}
     -k\left\lceil\frac{10}{\sqrt\eps}\right\rceil.
 \label{eq:subtract}
\end{equation}
Since $g\le4\eps$ and $\lceil10/\sqrt\eps\rceil\le11/\sqrt\eps$,
a fixed universal condition $\log n\ge44/c_0$ absorbs the overhead and gives
\[
 q\ge\frac{c_0}{4}\frac{k\log n}{\sqrt\eps}.
\]
All imported dimension conditions are polynomial in $k$ and $1/g$.
Together with the fixed threshold just used, they are implied by
$n\ge C_*(k/\eps)^D$ for universal $C_*,D$. This completes the proof of
Theorem~\ref{thm:large}.

\begin{remark}[What the reduction does and does not remove]
The postprocessor eliminates a possible $\log(1/\eps)$ overhead from a
cruder angle-amplification proof. It does not eliminate the dimension
hypothesis of the imported PCA construction. Nor does the subtraction in
\eqref{eq:subtract} by itself prove a useful bound in every small dimension.
\end{remark}

\section{Exact padding: transferring a hard lower rank}
\label{sec:padding}

\begin{lemma}[Rank-padding extraction without loss]
\label{lem:padding}
Let $p\ge0$, let $B\in\R^{m\times m}$, and set
$A=\diag(LI_p,B)$. Suppose $Z\in\R^{(p+m)\times(p+r)}$ has orthonormal
columns and $\norm{A(I-ZZ^T)}_2\le\tau<L$. Then an orthonormal
$U\in\R^{m\times r}$ can be computed from $Z$ without queries such that
\begin{equation}
 \norm{B(I-UU^T)}_2\le\tau.
 \label{eq:paddingconclusion}
\end{equation}
\end{lemma}
\begin{proof}
Write $Z=[C;D]$, where $C$ has $p$ rows. If $C$ did not have full row rank,
there would be a nonzero vector $(a,0)$ perpendicular to $\range Z$.
Its norm is multiplied by $L$ under $A$, contradicting $\tau<L$.
Thus $\dim\ker C=r$. Let $N$ have orthonormal columns spanning $\ker C$
and put $U=DN$. Then
$U^TU=N^TZ^TZN=I_r$.

For $x\perp\range U$, we have $N^TD^Tx=0$ and hence
$D^Tx\in\range C^T$. Choose $a$ with $C^Ta=-D^Tx$.
The vector $(a,x)$ is perpendicular to $\range Z$, so
\[
 L^2\norm a^2+\norm{Bx}^2
 \le\tau^2(\norm a^2+\norm x^2).
\]
Since $L>\tau$, this implies $\norm{Bx}\le\tau\norm x$.
Taking the supremum over $x\perp\range U$ proves the result.
The case $p=0$ is the same assertion with $U=Z$.
\end{proof}

\subsection{Application to the bounded rank-one hard family}
Let $m=n-k+1$, $p=k-1$, and $r=1$. On the hard rank-one family used for
\eqref{eq:paddedBN}, $\norm B_2<2$ and $\sigma_2(B)=1$. Use
$A=\diag(4I_{k-1},B)$. Each query to $A$ or $A^T$ can be simulated using
one query to the corresponding oracle for $B$; the other block is known.
Also $\sigma_{k+1}(A)=\sigma_2(B)=1$.

On every successful run, $\tau=1+\eps<4$, so
Lemma~\ref{lem:padding} extracts a valid rank-one output for $B$ at the
same accuracy and with no additional query. On a failure run the reduction
may return an arbitrary unit vector if extraction is undefined. Thus the
success guarantee is preserved, proving \eqref{eq:paddedBN} from the cited
rank-one lower bound. This argument uses a known norm bound on the hard
family; it does not assume that an arbitrary unbounded matrix can be scaled
for free without access to its norm.

\section{Why several tempting extensions are insufficient}
\label{sec:gaps}

\paragraph{Multiplying a rank-one lower bound by $k$.}
A direct sum of $k$ visible hard blocks is not a direct-sum theorem for this
oracle. A single query has components in all blocks and queries them all in
parallel. For example, if each visible block is a rank-one orthogonal
projector, one random product reveals a spanning vector of each block's
range simultaneously. A proof against all adaptive algorithms needs an
additional argument before a factor of $k$ is valid.

\paragraph{Confusing a valid residual with a good PCA objective.}
Example~\ref{ex:sharp} meets the spectral requirement while its captured
PCA fraction tends to zero. The weighted postprocessing theorem repairs this
specific obstruction; quoting a PCA theorem alone does not.

\paragraph{Confusing subspace depth with vector-query count.}
A lower bound on the number of block iterations is not automatically a
lower bound multiplied by a chosen block size for every adaptive algorithm.
Both the upper-bound implementation and the postprocessor above explicitly
count individual vector products.

\paragraph{Taking the minimum with $n$ on both sides.}
An exact-learning algorithm supplies an upper cap. It does not convert a
lower bound valid under $n\ge C_*(k/\eps)^D$ into a lower bound outside
that region. Doing so would discard a quantified hypothesis of the proof.

\subsection*{The remaining mathematical task}
A complete answer to RA-14 needs matching, universally quantified bounds in
the regimes not covered by Theorem~\ref{thm:large} or $k\ge n/2$.
For example, the current argument does not resolve arbitrary intermediate
powers of $k$ relative to $n$, nor accuracy regimes that violate the imported
dimension threshold. This note supplies no replacement hard distribution or
all-adaptive argument that closes those gaps. It also does not claim the
known global upper bound is the sharp expression in every such regime.

\section{Reproducible numerical checks}
\label{sec:experiments}

The package contains an oracle-counted implementation of the exact-column
method, promised-rank recovery, right block Krylov, and the symmetric Ritz
postprocessor. Dense residual norms and eigenvalues are computed only by
validation code, not inside the oracle algorithms. The mathematical model
uses exact real arithmetic; the code uses double precision.

All 18 unit tests passed in the supplied run. They cover query accounting,
rank-deficient cases, adaptive Gaussian completion identities, exact symbolic
polynomial identities and constants, the sharp counterexample, padding,
and the noncommuting weighted trace inequality. A finite set of tests does
not verify a universal statement over adaptive algorithms.

\subsection{Block Krylov trials, including observed failures}
For $n=96$, $k=4$, and $\eps=0.1$, each row below summarizes 30 trials with
independent random left and right orthogonal bases. The target residual
ratio is $1.1$. No theorem-certified universal depth constant is inferred
from these trials.

\begin{center}
\small
\begin{tabular}{@{}lrrrr@{}}
\toprule
Spectrum family & Depth & Queries & Successes & Worst ratio\\
\midrule
Flat spike & 2 & 24 & 12/30 & 1.199993\\
Flat spike & 4 & 40 & 29/30 & 1.199745\\
Flat spike & 8 & 72 & 30/30 & 1.000227\\
Slow decay & 2 & 24 & 30/30 & 1.089883\\
Slow decay & 4 & 40 & 30/30 & 1.020983\\
Slow decay & 8 & 72 & 30/30 & 1.000000\\
Large dynamic range & 2 & 24 & 30/30 & 1.001784\\
Large dynamic range & 4 & 40 & 30/30 & 1.000000\\
Large dynamic range & 8 & 72 & 30/30 & 1.000000\\
\bottomrule
\end{tabular}
\end{center}
The flat-spike family has four singular values equal to $1.2$ and a bulk
from $1$ down to $0.05$. Slow decay uses $0.96^j$. The third family starts
with $10^4,10^3,10,2$ and then decreases from $1$ to $0.01$.
All failures are retained in \code{results/krylov\_trials.csv}.

\subsection{Warm starts and scalar checks}
Sixty rotated instances were generated from the contraction representation
\eqref{eq:contract}, which ensures the starting residual bound. The
postprocessor met its trace target in all 60 trials; the smallest final
trace fraction was $0.9999999999999971$. No query budget was exceeded.
These small examples often converge much more accurately than the proved
bound, and some Krylov spaces fill the ambient space.

Twenty scalar grids, with $\eps$ from $10^{-8}$ to $0.49$, tested
\eqref{eq:coefficient}. The largest observed coefficient divided by $\eps$
was approximately $3.472\cdot10^{-10}$, below the proved $1/200$ bound.
The grid is only a numerical check; Section~\ref{sec:filter} supplies the
continuous-domain proof. The maximum orthogonality error across the 270
block-Krylov trials was $6.242\cdot10^{-15}$.

\subsection{Commands and limitations}
From the extracted package directory:
\begin{verbatim}
python -m pip install -r requirements.txt
OPENBLAS_NUM_THREADS=1 python -m unittest discover -s tests -v
OPENBLAS_NUM_THREADS=1 python run_experiments.py
pdflatex -interaction=nonstopmode -halt-on-error report.tex
pdflatex -interaction=nonstopmode -halt-on-error report.tex
\end{verbatim}
The recorded environment was Python 3.13.5, NumPy 2.3.5, and SymPy 1.14.0.
Floating-point results can vary slightly with the numerical library. There
is no formal verification in a proof assistant. The imported asymptotic
constants are not estimated by the experiments.

\section{Conclusion}
The note proves the global rank lower bound and the $\Theta(n)$ regime
$k\ge n/2$, and develops an exact padding transfer. Its main additional
argument is a spectral-to-PCA postprocessor costing
$k\lceil10/\sqrt\eps\rceil$ products. Combined with a published adaptive
PCA lower bound, this gives matching $\Theta(k\log n/\sqrt\eps)$ bounds
under a polynomial large-dimension hypothesis, with growing $k$ allowed.

\textbf{This does not fully solve RA-14.} The unresolved part of this
write-up is the simultaneous finite-parameter characterization outside
those regimes. The source files, proof dependencies, numerical checks, and
this limitation are included together so that the established reductions
can be inspected without mistaking them for a complete answer.

\begingroup
\footnotesize
\raggedright
\setlength{\parskip}{0pt}
\begin{thebibliography}{9}
\setlength{\itemsep}{2pt}
\bibitem{RA14}
Open Problems in Numerical Linear Algebra.
\emph{RA-14: Optimal query complexity of spectral rank-$k$ approximation}.
Repository problem statement, accessed September 12, 2026.
\url{https://github.com/ajt60gaibb/OpenProblemsInNLA/tree/main/randomized-and-low-rank-approximation/RA-14}.

\bibitem{MM15}
Cameron Musco and Christopher Musco.
\emph{Randomized Block Krylov Methods for Stronger and Faster Approximate
Singular Value Decomposition}. 2015. arXiv:1504.05477.
\url{https://arxiv.org/abs/1504.05477}.

\bibitem{SAR18}
Max Simchowitz, Ahmed El Alaoui, and Benjamin Recht.
\emph{Tight Query Complexity Lower Bounds for PCA via Finite Sample
Deformed Wigner Law}. STOC 2018; arXiv:1804.01221v2, 2020.
\url{https://arxiv.org/abs/1804.01221}.

\bibitem{BN23}
Ainesh Bakshi and Shyam Narayanan.
\emph{Krylov Methods are (nearly) Optimal for Low-Rank Approximation}.
2023. arXiv:2304.03191v1.
\url{https://arxiv.org/abs/2304.03191}.
\end{thebibliography}
\endgroup
\end{document}
